% Generated draft: requires mathematical and source review.
\documentclass[11pt]{article}
\usepackage[margin=25mm]{geometry}
\usepackage{fontspec}
\setmainfont{Latin Modern Roman}
\usepackage{amsmath,amssymb,mathtools,unicode-math}
\setmathfont{Latin Modern Math}
\usepackage{enumitem,needspace,xurl,hyperref}
\setlength{\parindent}{0pt}
\setlength{\parskip}{5pt}
\setlength{\emergencystretch}{3em}
\newcommand{\N}{\mathbb{N}}
\newcommand{\Z}{\mathbb{Z}}
\newcommand{\Q}{\mathbb{Q}}
\newcommand{\R}{\mathbb{R}}
\newcommand{\C}{\mathbb{C}}
\newcommand{\F}{\mathbb{F}}
\newcommand{\A}{\mathbb{A}}
\newcommand{\PP}{\mathbb{P}}
\newcommand{\E}{\mathbb{E}}
\newcommand{\eps}{\varepsilon}
\newcommand{\dd}{\,\mathrm{d}}
\begin{document}
\section*{UnsolvedMath ID: 800308}
\textbf{Source problem code:} catalogue entry
\subsection{Definitions and mathematical statement}
Let \(\mathbb T=\R/\Z\), and let \(\operatorname{Diff}^{2}(\mathbb T)\) denote the group of circle diffeomorphisms whose maps and inverses are of class \(C^{2}\). For an orientation-preserving \(f\in\operatorname{Diff}^{2}(\mathbb T)\), its rotation number is
\[
\rho(f)=\left[\lim_{n\to\infty}\frac{F^{n}(x)-x}{n}\right]\in\R/\Z,
\]
where \(F:\R\to\R\) is a lift satisfying \(F(x+1)=F(x)+1\); the class is independent of the lift and of \(x\). For \(\rho\in\R/\Z\), write \(R_{\rho}([x])=[x]+\rho\). An irrational rotation number means \(\rho\notin\Q/\Z\).

Define the set of \(C^{2}\) conjugates by
\[
\mathcal C^{2}(f)=\{h\circ f\circ h^{-1}:h\in\operatorname{Diff}^{2}(\mathbb T)\}.
\]
The question is whether, for every orientation-preserving \(C^{2}\) circle diffeomorphism \(f\) with irrational rotation number \(\rho\),
\[
R_{\rho}\in\overline{\mathcal C^{2}(f)}^{\,C^{2}}.
\]
Explicitly, fixing a real representative \(\alpha\) of \(\rho\), does there exist a sequence \(h_n\in\operatorname{Diff}^{2}(\mathbb T)\) and lifts \(G_n\) of \(h_n\circ f\circ h_n^{-1}\) such that
\[
\lim_{n\to\infty}\max_{0\leq j\leq 2}\sup_{x\in[0,1]}
\left|\frac{\mathrm d^{j}}{\mathrm dx^{j}}\bigl(G_n(x)-(x+\alpha)\bigr)\right|=0?
\]
Here the closure uses uniform convergence of lifts and their derivatives through order two on a fundamental interval, with compatible choices of lifts. The supplied source does not explicitly specify whether conjugating diffeomorphisms must preserve orientation; the displayed conjugacy set allows both orientations. This convention requires confirmation against the intended source formulation.
\subsection{Short English statement}
Can every \(C^{2}\) circle diffeomorphism with irrational rotation number \(\rho\) be conjugated by \(C^{2}\) diffeomorphisms so that the resulting maps converge, together with their first two derivatives, to the rigid rotation \(R_{\rho}\)? The supplied registry labels the problem open; this draft does not independently verify that status.
\subsection{Sources}
\begin{itemize}
\item Local UnsolvedMath registry, supplied record with permanent ID \(800308\), source problem code \texttt{catalogue entry}, in the AMR Open Problem Lists collection. The record attributes the question to Navas, with proposed year \(2017\), and supplies the status \emph{open}.
\item Navas, \emph{Group Actions on 1-Manifolds: A List of Very Concrete Open Questions} (2017), \url{https://arxiv.org/abs/1712.06462}. The supplied title identifies Question 25, and the provenance identifies the source item as ``Question 25 (\#24).'' The relationship between these two numbers is not explained in the record.
\item Supplied provenance records access on \(2026\text{-}10\text{-}16\) and extraction by \texttt{arxiv\_latex}. These are supplied provenance details, not a claim of independent source access or verification. The supplied upstream revision is \texttt{d371ee917ac9647fbb06cd6fc529deec6f3dda82}.
\end{itemize}
\end{document}
